\begin{figure}[H]
\centering
\resizebox{0.95\textwidth}{!}{%
\begin{tikzpicture}[
    >=stealth, font=\sffamily,
    % Styles
    threadBox/.style={draw, thick, rounded corners=8pt, inner sep=15pt},
    actionNode/.style={draw, thick, fill=white, rounded corners=4pt, minimum width=2.8cm, minimum height=0.8cm, align=center},
    queueNode/.style={draw, thick, rounded corners=10pt, minimum width=3.5cm, minimum height=1cm, align=center},
    dbNode/.style={draw, thick, fill=gray!15, cylinder, shape border rotate=90, aspect=0.25, minimum width=1.4cm, minimum height=1.2cm, align=center}
]

% ================= DATABASES (Bờ ngoài cùng) =================
\node[dbNode] (minio) at (0, 4) {\textbf{MinIO}};
\node[dbNode] (neo) at (5, 0.0) {\textbf{Neo4j}};

% ================= CỘT 1: MAIN THREAD (Non-blocking) =================
\node[actionNode, draw=blue!80] (pdf) at (0, 6.5) {\textbf{pdf extractor}\\};
\node[actionNode, dashed, draw=gray, text=gray] (other) at (0, 1.5) {\textit{Main Thread Free}\\\textit{(Repeat PDF extraction...)}};

% ================= CỘT 2: MIDDLEWARE QUEUES (Buffer) =================
\node[queueNode, fill=purple!10, draw=purple!60] (q1) at (5, 6.5) {\texttt{extract\_queue}\\(Chunk Buffer)};
\node[queueNode, fill=green!10, draw=green!60] (q2) at (5, 4) {\texttt{build\_queue}\\(Graph Buffer)};

% ================= CỘT 3: ASYNC WORKER POOL =================
% 1. Parallel LLM Workers (Vẽ 3 hộp đè lên nhau tạo hiệu ứng Parallel)
\node[actionNode, fill=orange!10, draw=orange!80] (w3) at (10.6, 6.9) {\textbf{LLM Worker 2}};
\node[actionNode, fill=orange!10, draw=orange!80] (w2) at (10.8, 6.7) {\textbf{LLM Worker 1}};
\node[actionNode, fill=orange!10, draw=orange!80] (w1) at (11.0, 6.5) {\textbf{LLM Worker 0}};

% 2. Sequential Graph Constructor
\node[actionNode, fill=teal!10, draw=teal!80] (gc) at (5, 2.25) {\textbf{Graph Constructor}\\(Sequential)};


% ================= TIMELINE (Trục dọc) =================
\draw[->, thick, gray] (13, 7.5) -- (13, 0);
\node[right, font=\small] at (13, 6.5) {Time $t_0$};
\node[right, font=\small] at (13, 4.75) {Time $t_1$};
\node[right, font=\small] at (13, 3) {...............};
\node[right, font=\small] at (13, 1) {Time $t_n$};

% ================= ARROWS & ROUTING =================
% Main Thread activities
\draw[->, thick, blue!80] (pdf.south) -- (minio.north) node[midway, above, font=\scriptsize] {upload\_raw};
\draw[->, thick, blue!80] (pdf.east) -- (q1.west) node[midway, above, font=\scriptsize] {put(chunk)};
\draw[->, thick, dashed, gray] (minio.south) -- (other.north) node[midway, right, font=\scriptsize, align=left] {Keep running, No waiting};

% Queue 1 to Workers
\draw[->, thick, orange!80] (q1.east) -- (w1.west) node[midway, above, font=\scriptsize] {get()};

% Workers to Queue 2 (Đường chéo vắt ngược lại Buffer)
\draw[->, thick, orange!80] (w1.south) -- (q2.north east) node[midway, above right, font=\scriptsize, xshift=0.3cm] {put(sub-graph)};

% Queue 2 to Graph Constructor
\draw[->, thick, teal!80] (q2.south) -- (gc.north) node[midway, above, font=\scriptsize] {get()};

% Graph Constructor to Neo4j
\draw[->, thick, teal!80] (gc.south) -- (neo.north) node[midway, above, font=\scriptsize] {update()};

\end{tikzpicture}%
}
\vspace{15pt}\caption{Asynchronous Ingestion Pipeline integrating Queues, Parallel Workers, and Storage}
\label{fig:ingestion}
\end{figure}